% TOP_PROBLEM: {"id": 30005276, "title": "Density Hypothesis for Riemann-Zeta Zeros", "category_id": 1, "category": {"id": 1, "name": "number_theory", "display_name": "Number Theory", "description": "Properties of integers, prime numbers, Diophantine equations.", "slug": "number-theory", "order_index": 1, "created_at": "2026-07-31T15:26:25.670Z"}, "status": "partially_solved"}
% FIRST_POSTED: 2026-09-25
% ENTRY_KIND: statement_only
% !TeX program = lualatex
% Definition and sources only; this file is not an LLM solution attempt.
\documentclass[11pt]{article}
\usepackage[margin=25mm]{geometry}
\usepackage{fontspec}
\setmainfont{Latin Modern Roman}
\usepackage{amsmath,amssymb,mathtools}
\usepackage{unicode-math}
\setmathfont{Latin Modern Math}
\usepackage{xurl,hyperref}
\hypersetup{colorlinks=true,urlcolor=blue}
\setcounter{secnumdepth}{0}
\setlength{\parindent}{0pt}
\setlength{\parskip}{5pt}
\setlength{\emergencystretch}{3em}
\begin{document}
\section*{\texorpdfstring{Density Hypothesis for Riemann-Zeta Zeros}{Density Hypothesis for Riemann-Zeta Zeros}}\label{30005276}
\subsection{Definitions and mathematical statement}
Count zeta zeros with multiplicity by
\[
 N(\sigma,T)=\#\{\rho=\beta+i\gamma:\zeta(\rho)=0,\ 0<\beta<1,
                                  \beta\ge\sigma,\ |\gamma|\le T\}.
\]
For each fixed $1/2\le\sigma\le1$ and $\varepsilon>0$, the density hypothesis predicts
\[
 N(\sigma,T)\ll_{\sigma,\varepsilon}T^{2(1-\sigma)+\varepsilon}.
\]
The bound is uniform in large $T$, with dependence on $\sigma,\varepsilon$ permitted. It limits the number of off-line zeros rather than asserting that they do not exist.
\par\smallskip\textit{Formulation and concept references:} \href{https://terrytao.wordpress.com/tag/riemann-zeta-function/}{[S1]}.

\subsection{Short English statement}
Are zeta zeros to the right of any fixed vertical line as sparse as the density hypothesis predicts?
\subsection{Sources}
\begin{itemize}
\item[S1]Terence Tao, "A computation-outsourced discussion of zero density theorems for the Riemann zeta function", What's new, 7 July 2024.
\url{https://terrytao.wordpress.com/tag/riemann-zeta-function/}.
\textit{Formulation source.}
\end{itemize}
\end{document}
